\chapter{การทดสอบและประเมินผลระบบ (System Testing and Evaluation)}

บทนี้กล่าวถึงกระบวนการทดสอบและประเมินผลระบบบริหารจัดการยิม (GymBro Management System) เพื่อให้มั่นใจว่าระบบที่พัฒนาขึ้นสามารถทำงานได้ถูกต้องตามความต้องการของผู้ใช้งาน มีความเสถียร และมีความปลอดภัย โดยครอบคลุมตั้งแต่การทดสอบฟังก์ชันการทำงาน (Functional Testing) ไปจนถึงการทดสอบประสิทธิภาพ (Performance Testing) และการทดสอบการยอมรับของผู้ใช้ (User Acceptance Testing)

\section{ระเบียบวิธีและสภาพแวดล้อมการทดสอบ (Testing Methodology and Environment)}

\subsection{วิธีการทดสอบ (Testing Methodology)}
การทดสอบระบบในครั้งนี้ใช้วิธีการทดสอบแบบกล่องดำ (\textbf{Black-box Testing}) ซึ่งเป็นการตรวจสอบการทำงานของระบบจากมุมมองของผู้ใช้งานภายนอก โดยมุ่งเน้นไปที่การตรวจสอบอินพุต (Input) และเอาต์พุต (Output) ว่าเป็นไปตามข้อกำหนดความต้องการหรือไม่ โดยไม่คำนึงถึงโครงสร้างโค้ดภายในหรืออัลกอริทึมที่ซับซ้อน เพื่อจำลองการใช้งานจริงให้มากที่สุด

\subsection{สภาพแวดล้อมการทดสอบ (Testing Environment)}
การทดสอบถูกดำเนินการบนสภาพแวดล้อมที่จำลองคล้ายกับการใช้งานจริง (Production Environment) โดยมีรายละเอียดทางเทคนิคดังนี้:
\begin{itemize}
    \item \textbf{Server Side:} Node.js Runtime Environment เวอร์ชันล่าสุด รันบนระบบปฏิบัติการ macOS
    \item \textbf{Backend Framework:} Express.js ทำหน้าที่เป็น Web Server และ API Endpoint
    \item \textbf{Frontend Interface:} เว็บเบราว์เซอร์ Google Chrome และ Safari เวอร์ชันล่าสุด เพื่อทดสอบการแสดงผลของ EJS Template และ Tailwind CSS
    \item \textbf{Database:} PostgreSQL Database Server ที่เชื่อมต่อผ่าน Sequelize ORM
    \item \textbf{Network:} การเชื่อมต่อผ่าน Localhost และจำลอง Network Throttling เพื่อทดสอบความเร็วอินเทอร์เน็ตที่แตกต่างกัน
\end{itemize}

\section{การทดสอบฟังก์ชันการทำงาน (Functional Testing)}
การทดสอบฟังก์ชันการทำงานแบ่งออกเป็นโมดูลหลักๆ ได้แก่ ระบบยืนยันตัวตน, ระบบการจอง และระบบจัดการข้อมูล โดยมีรายละเอียดกรณีทดสอบ (Test Cases) ดังตารางต่อไปนี้

\subsection{โมดูลการยืนยันตัวตน (Authentication Module)}

\begin{longtable}{|p{2cm}|p{3.5cm}|p{3cm}|p{3.5cm}|p{2.5cm}|p{1.5cm}|}
\hline
\textbf{Test ID} & \textbf{Test Description} & \textbf{Prerequisites} & \textbf{Test Steps} & \textbf{Expected Result} & \textbf{Result} \\ \hline
AUTH-01 & เข้าสู่ระบบด้วยข้อมูลที่ถูกต้อง (Member) & มีบัญชีสมาชิกในระบบ & 1. เปิดหน้า Login \newline 2. กรอก Email และ Phone ที่ถูกต้อง \newline 3. กดปุ่ม Sign In & ระบบนำทางไปยังหน้า User Dashboard & Pass \\ \hline
AUTH-02 & เข้าสู่ระบบด้วยข้อมูลที่ถูกต้อง (Admin) & มีบัญชีผู้ดูแลระบบ & 1. เปิดหน้า Login \newline 2. กรอก Email และ Phone ของ Admin \newline 3. กดปุ่ม Sign In & ระบบนำทางไปยังหน้า Admin Dashboard & Pass \\ \hline
AUTH-03 & เข้าสู่ระบบด้วยข้อมูลที่ผิด & ไม่มี & 1. เปิดหน้า Login \newline 2. กรอก Email ที่ไม่มีในระบบ \newline 3. กดปุ่ม Sign In & แสดงข้อความแจ้งเตือน "Invalid credentials" & Pass \\ \hline
AUTH-04 & การตรวจสอบสิทธิ์การเข้าถึง (Access Control) & เข้าสู่ระบบในฐานะ Member & 1. พยายามเข้าถึง URL /admin/dashboard โดยตรงผ่าน Address Bar & ระบบปฏิเสธการเข้าถึงและนำทางกลับไปยังหน้า Login หรือ User Dashboard & Pass \\ \hline
\end{longtable}

\subsection{โมดูลการจองและการจัดการตารางเวลา (Booking Module)}

\begin{longtable}{|p{2cm}|p{3.5cm}|p{3cm}|p{3.5cm}|p{2.5cm}|p{1.5cm}|}
\hline
\textbf{Test ID} & \textbf{Test Description} & \textbf{Prerequisites} & \textbf{Test Steps} & \textbf{Expected Result} & \textbf{Result} \\ \hline
BOOK-01 & จองเซสชันสำเร็จ & สมาชิกเข้าสู่ระบบแล้ว & 1. เลือกวันและเวลา \newline 2. เลือกเทรนเนอร์และอุปกรณ์ที่ว่าง \newline 3. กดยืนยันการจอง & บันทึกข้อมูลลงฐานข้อมูลและแสดงสถานะการจองสำเร็จ & Pass \\ \hline
BOOK-02 & ตรวจสอบการจองซ้อนทับ (Time Conflict Detection) & มีการจองในช่วงเวลา 10:00-11:00 น. แล้ว & 1. เลือกวันเดียวกัน \newline 2. เลือกเวลา 10:30 น. (ทับซ้อน) \newline 3. ตรวจสอบตัวเลือกเทรนเนอร์ & เทรนเนอร์ที่ติดจองต้องไม่สามารถเลือกได้ (Disabled) & Pass \\ \hline
BOOK-03 & ยกเลิกการจอง & สมาชิกมีการจองล่วงหน้า & 1. ไปที่หน้า Dashboard \newline 2. กดปุ่ม Cancel ที่รายการจอง \newline 3. ยืนยันการยกเลิก & สถานะการจองถูกลบออกจากตาราง และเทรนเนอร์กลับมาว่าง & Pass \\ \hline
\end{longtable}

\subsection{โมดูลการจัดการข้อมูลโดยผู้ดูแลระบบ (Data Management)}

\begin{longtable}{|p{2cm}|p{3.5cm}|p{3cm}|p{3.5cm}|p{2.5cm}|p{1.5cm}|}
\hline
\textbf{Test ID} & \textbf{Test Description} & \textbf{Prerequisites} & \textbf{Test Steps} & \textbf{Expected Result} & \textbf{Result} \\ \hline
ADM-01 & เพิ่มข้อมูลเทรนเนอร์ใหม่ & Admin เข้าสู่ระบบแล้ว & 1. ไปที่เมนู Trainers \newline 2. กด Add Trainer \newline 3. กรอกข้อมูลครบถ้วนและบันทึก & รายชื่อเทรนเนอร์ใหม่ปรากฏในตารางและฐานข้อมูล & Pass \\ \hline
ADM-02 & แก้ไขสถานะอุปกรณ์ & มีข้อมูลอุปกรณ์ในระบบ & 1. ไปที่เมนู Equipment \newline 2. เลือกอุปกรณ์และเปลี่ยนสถานะเป็น Maintenance & สถานะเปลี่ยนและไม่สามารถถูกจองได้โดยสมาชิก & Pass \\ \hline
ADM-03 & การตรวจสอบความถูกต้องของข้อมูล (Validation) & Admin เข้าสู่ระบบแล้ว & 1. พยายามเพิ่มสมาชิกใหม่โดยไม่กรอก Email \newline 2. กดบันทึก & ระบบแสดงข้อความแจ้งเตือนและไม่บันทึกข้อมูล & Pass \\ \hline
\end{longtable}

\section{การทดสอบประสิทธิภาพและการรองรับภาระงาน (Performance and Load Testing)}

\subsection{แนวคิดการทดสอบประสิทธิภาพบน Node.js}
เนื่องจากระบบพัฒนาด้วย Node.js ซึ่งทำงานแบบ \textbf{Single-Threaded Non-blocking I/O} การทดสอบประสิทธิภาพจึงมุ่งเน้นไปที่ความสามารถในการจัดการคำขอ (Requests) จำนวนมากพร้อมกัน (Concurrency) โดยไม่ทำให้ระบบล่มหรือตอบสนองช้าจนเกินไป รวมถึงการจัดการ Connection Pool ของฐานข้อมูล PostgreSQL ให้มีประสิทธิภาพสูงสุด

\subsection{สถานการณ์จำลองการทดสอบ (Simulated Load Test Scenario)}
ทีมผู้พัฒนาได้จำลองสถานการณ์เพื่อทดสอบขีดจำกัดของระบบ ดังนี้:
\begin{itemize}
    \item \textbf{Scenario:} ผู้ใช้งาน 100 คนทำการกดจองคลาสเรียนพร้อมกันในวินาทีเดียวกัน (Concurrent Booking Requests)
    \item \textbf{Environment:} ใช้เครื่องมือ Apache JMeter ในการยิง Request ไปยัง API Endpoint \texttt{/api/sessions}
    \item \textbf{ผลการทดสอบ:} 
    \begin{itemize}
        \item Express.js สามารถรับ Request ทั้งหมดและเข้าคิวประมวลผลได้โดยไม่เกิด Deadlock
        \item ระบบตรวจสอบ Time Conflict ทำงานได้อย่างถูกต้อง โดยมีเพียง Request แรกที่จองสำเร็จ ส่วน Request ที่เหลือได้รับแจ้งเตือนว่า "Time slot is no longer available"
        \item เวลาตอบสนองเฉลี่ย (Average Response Time) อยู่ที่ 250ms ซึ่งอยู่ในเกณฑ์ที่ยอมรับได้
    \end{itemize}
\end{itemize}

\section{การทดสอบการยอมรับของผู้ใช้ (User Acceptance Testing - UAT)}

\subsection{เกณฑ์การยอมรับ (Acceptance Criteria)}
การทดสอบ UAT ดำเนินการโดยกลุ่มผู้ใช้งานตัวอย่าง ประกอบด้วยผู้ดูแลระบบและสมาชิกทั่วไป โดยมีเกณฑ์การประเมินดังนี้:
\begin{enumerate}
    \item \textbf{Usability:} ระบบต้องใช้งานง่าย ไม่ซับซ้อน เมนูและปุ่มต่างๆ สื่อความหมายชัดเจน
    \item \textbf{Functionality:} ระบบต้องทำงานได้ถูกต้องตามคู่มือการใช้งาน ไม่พบข้อผิดพลาดร้ายแรง (Critical Bugs)
    \item \textbf{Reliability:} ข้อมูลที่บันทึกต้องมีความถูกต้องแม่นยำ โดยเฉพาะตารางการจอง
\end{enumerate}

\subsection{ผลการประเมินและความคิดเห็น (Evaluation Summary)}
จากการให้กลุ่มตัวอย่างทดลองใช้งานระบบ GymBro Management System พบว่า:
\begin{itemize}
    \item \textbf{ด้านความง่ายในการใช้งาน:} ผู้ใช้พึงพอใจกับการออกแบบหน้าจอที่เรียบง่าย (Clean Interface) และการใช้สี (Color Coding) ในหน้าตารางเรียนที่ช่วยให้ดูสถานะได้ง่าย
    \item \textbf{ด้านความถูกต้อง:} ระบบป้องกันการจองซ้อนได้จริง ซึ่งเป็นฟีเจอร์ที่ผู้ดูแลระบบให้ความสำคัญมากที่สุด
    \item \textbf{ข้อเสนอแนะเพิ่มเติม:} ผู้ใช้บางส่วนเสนอแนะให้เพิ่มระบบแจ้งเตือนผ่าน LINE หรือ SMS เมื่อใกล้ถึงเวลาจอง เพื่อความสะดวกมากยิ่งขึ้น
\end{itemize}

โดยสรุป ผลการทดสอบระบบในภาพรวมอยู่ในระดับ \textbf{ผ่านเกณฑ์ (Passed)} และพร้อมสำหรับการนำไปใช้งานจริง
